\section{朝廷、法律、家庭、城市与海外网络}
\label{sec:ming-court-law-family-city-overseas}

明代政治不能仅由皇帝、内阁与党争概括。律例、审判、婚姻、城市市场、书院寺庙、海港和海外往来，使朝廷的决定经由地方中介进入不同家庭。实录、会典和奏疏能钉住任命、法律与重大政争；判牍、契约、墓志、家谱、方志、笔记、外来旅行者与贸易材料各自只显示局部实践。它们不容许把一条法律、一座城市或一艘商船转化为整个帝国的社会平均数。

\subsection{宫廷信息与地方执行}

洪武至永乐的皇权、迁都、军政与法令形成新的中央制度，嘉靖、万历、天启与崇祯时期的奏疏、廷争和人事则显示信息渠道不断紧张。朝廷发布的命令须经过部院、督抚、州县、书吏与地方精英才可能落实；上谕存在不等于地方已经理解或执行。将党争完全解释为个人道德，或将其视为与饷银、灾害和地方征解无关的纯宫廷戏剧，都过于简化。

内阁票拟、司礼监、科道言官和地方督抚构成不同的文书与责任链。奏折具有请饷、辩护、攻击或邀功的目的，不能把其数字和危机陈述当作透明现实。尤其晚明军费、赈济、矿税与海防争论，应分辨定额、解送、实际收到和地方家庭所承受的不同层次。

\subsection{律例、诉讼与性别化家庭}

《大明律》及后续条例规定户婚、田土、债务、刑罚和官民关系，是研究国家规范的重要入口。法律条文说明应然秩序，不能证明所有人均能进入衙门、获得同样证据或面对同样裁判。路程、费用、识字、宗族、身份和地方官能力都会影响诉讼。判牍保存的是进入案件的冲突，不能直接代表平日所有家庭关系。

婚姻、嫁妆、继承、收养、寡居与妇女劳动常在契约、家训、节孝叙述或诉讼中断续可见。家训和旌表特别容易呈现精英希望规范的贞节与孝道，不能把规范语言等同于家庭实际。妇女的财产、经营、迁徙和宗教活动不应被抹去，也不应由少数有产、能留名者代表城乡和族群差异极大的女性经验。

\subsection{城市文化与公共社群}

南京、北京、苏州、杭州、扬州、泉州、广州及众多府城的市场、书坊、园林、会馆、寺观和戏曲空间，使城市文化与区域贸易相连。文人笔记和画作擅长保存消费、交游和审美，容易偏向富有男性；税收、价格、工匠和雇工的材料则不连续。城市繁华不能自动说明周边农村富裕，也不能将贫困叙述扩展为所有市民的生活。

书院、社学、宗祠、寺院、善会和会馆既可传播知识，也能组织慈善、借贷、旅居互助和地方声望。碑刻与捐册说明部分捐施者和仪式，不能量化参与者，更不能将宗教、商帮或士绅写成内部无分歧的实体。城市公共生活常在国家规制、家族网络和市场利益之间协调。

\subsection{海禁、海防与海外连接}

朝廷的海禁、倭患、互市和朝贡政策，属于安全、财政和外交的多重安排。禁令能证明国家忧虑和管理目标，不能直接证明沿海贸易停止；发现私贸、沉船或外来货物也不能证明法律无效。福建、广东、浙江和东南岛屿的港口、商人、渔民、军队与地方官面对不同的季风、战争和市场条件。

郑和下西洋、葡萄牙人在澳门活动、马尼拉贸易以及晚明传教士和白银流通，使中国与印度洋、东南亚、美洲和欧洲的网络相遇。航海记录、传教书信、海关材料与地方文献各有目的；它们可说明某些人、物和知识进入特定港口，不能把“全球化”写成所有地区同步的经验。白银的流入、赋役折纳、矿业和家庭借贷之间的关系也须分区、分时讨论。

\subsection{危机中的法律与城市秩序}

十七世纪的辽东战争、灾害、饷银和流民使京师、府城和县镇的治安、司法与慈善承压。崇祯末年政权崩解的年序能够确证，城市是否同样缺粮、地方是否同样失序则不能由一份奏报或一座城的见闻概括。法律、家庭、市场和海外联系在危机中并未消失，而是以不同地区、不同群体可见或不可见的方式被重新安排。
